\section{Discussion and limitations}\label{sec:discussion}

\paragraph{What changed the conclusions, and what did not.}
On this family, none of the conventions we tested changed which method was
cheapest. The accuracy target, the instance regime and the problem dimension
did change the order below it: the change of lower-place ordering with the
accuracy target that we found at $n=100$ does not reproduce at $n=500$
(\cref{sec:n500}), and the rank effect that organizes it at $n=100$ no longer
changes the order there. The stopping rule
changed the achieved accuracy by a factor of 383. The cost unit (\EVDs\ against wall-clock time) and the accounting rule
(accepted-only against all-trial) changed cost ratios but no ordering. Two
line-search details changed whether methods converged at all. We think the
correct summary is conservative. The protocol matters mainly for the
\emph{magnitude} of reported advantages and for the ordering of methods that
are close in cost. On the evidence here, it does not matter for identifying a
clear winner. We expect this to hold wherever one method is ahead by a
large factor; where methods are within a factor of about two, conclusions
drawn under a single convention should be treated with caution.

\paragraph{Why an exact family.}
The regime dependence in \cref{sec:cells} is organized mainly by rank, one of
the three quantities the construction controls. Finding it requires both an
exact $\Xs$, so that forward error is not contaminated by an approximate
reference, and independent control of the spectral structure at the
solution. Real matrices supply neither. The same machinery then showed the
limit of that finding: at $n=500$ (\cref{sec:n500}) the rank effect keeps its
direction but no longer changes the order, so what looked like a rule about
the methods at one dimension turns out to depend on the benchmark regime as
well.

\paragraph{Limitations.}
\begin{itemize}
  \item \emph{Problem size.} Results are reported at $n=100$ and $n=500$
    (270 instances each). The $r\ge20$ crossover between SBB-Dual and
    AGD-SDAJ found at $n=100$ does \emph{not} persist at $n=500$, where
    AGD-SDAJ-BH is cheaper at every rank. No dimension above 500 was tested
    on the family. Dykstra-APM is censored at $n=500$ by its iteration cap.
  \item \emph{Synthetic family.} The KKT family is synthetic. Its $G$
    matrices are correlation-like but are not estimated from data.
    \Cref{sec:real} adds four matrices from the literature, from $n=94$ to
    $n=3250$, measured against a computed rather than an exact reference; for
    the two largest that reference is capped and reaches
    $\nrm{\nabla\theta}_2\approx10^{-12}$ rather than $10^{-13}$, and no
    wall-clock times were taken. Four matrices cannot separate dimension from
    the other ways in which they differ.
  \item \emph{Statistical uncertainty.} Aggregate medians at both dimensions
    carry cluster-bootstrap intervals (\cref{tab:fwdci,tab:n500}). Each
    dimension rests on five paired seeds per rank, that is 15 seed clusters,
    so inferential statements remain modest. Cell-level
    orderings rest on five instances each, and within a cell the adjacent pair
    AGD-SDAJ-BH against AGD-SDAJ is usually unresolved, so the count of
    distinct orderings should not be read as a count of real regime changes.
    We report counts and intervals rather than significance tests, because the
    reuse of each drawn $U$ across all cells of its rank makes instance-level
    comparisons strongly dependent.
  \item \emph{Single machine.} All timings come from one machine with
    single-threaded BLAS.
  \item \emph{Our adaptation.} The Borsdorf--Higham adaptation to AGD-SDAJ is
    ours, not the original authors'. It is reported beside the unmodified
    method throughout.
  \item \emph{Feasibility propositions.} \Cref{prop:det} assumes uniform
    $\mu$, and the extension to mixed $\mu$ is open (\cref{app:family}).
    \Cref{prop:hd} is stated for $a$ only; no companion bound on $b$ is
    given, since the released family lies in the fixed-rank regime where
    \cref{prop:hd} does not apply.
  \item \emph{Scope of methods.} We benchmark five variants of four methods.
    Other methods (e.g.\ Anderson-accelerated alternating
    projections~\cite{higham2016anderson}) are not included.
    Neither are the recent semismooth Newton method for general projection
    equations~\cite{armijo2025semismooth} or the factor-structured
    variant~\cite{borsdorf2010factor}.
  \item \emph{Attribution of the line search.} We use the Borsdorf--Higham
    globalization but not the rest of their algorithm, and the adaptation of
    their near-equality safeguard to AGD-SDAJ is ours
    (\cref{sec:solvers,sec:trap2}).
  \item \emph{Unweighted problem.} Only the unweighted Frobenius problem is
    considered.
\end{itemize}

\paragraph{Conclusion.}
An exact instance family with a prescribed solution structure makes it
possible to separate stopping-rule, reference, accounting and
primitive-work effects in NCM benchmarks. On 270 instances at $n=100$, these
effects changed cost ratios and the lower places of the ranking, but not the
cheapest method. Repeating the study on 270 instances at $n=500$ sharpened
that finding: Newton-SIN-BH remains the leading method, but the change of
lower-place ordering with the accuracy target does not reproduce, and the rank
effect that organizes it at $n=100$ no longer changes the order. Rankings
below the leading method are therefore not intrinsic properties of the
solvers; they depend materially on the benchmark regime, including its
dimension. The one clear failure we found is a known finite-precision
line-search defect, which appears in two separately published methods, on
the constructed family and on real matrices. It is
removed in the semismooth Newton method by the published Borsdorf--Higham
modified line search, and in AGD-SDAJ by our adaptation of their near-equality
safeguard. We recommend
that NCM comparisons report forward error against a trusted solution at a
common tolerance, give rejected trials no accuracy credit, and report
primitive work and wall-clock time together.
